\section{Circuit Complexity --- Lupanov}
%%% generated by the blueprint pipeline from TCSlib.Complexity.CircuitComplexity.Lupanov
\subsection*{Circuit Complexity --- Lupanov: Introduction}

Declarations of the Lean module \texttt{TCSlib.Complexity.CircuitComplexity.Lupanov}.

\subsection*{Circuit Complexity --- Lupanov: Results}

\begin{theorem}[Reconstruction from heap paths]
\lean{BoolCircuit.Lupanov.assemble_heapPath}
\label{BoolCircuit.Lupanov.assemble_heapPath}
\leanok
\uses{BoolCircuit.Lupanov.assemble, BoolCircuit.Lupanov.heapPath, BoolCircuit.Lupanov.heapPath_bounds, BoolCircuit.Lupanov.testBit_heapPath, BoolCircuit.Lupanov.xget}
\difficulty{4}
Let $n,k\in\mathbb{N}$ with $k\le n$, and let $x\in\{0,1\}^n$. Assembling the two
heap-path indices of $x$, computed on the first $n-k$ and last $k$ coordinates
respectively, recovers $x$.
\end{theorem}

\begin{theorem}[Circuit output from combination roots]
\lean{BoolCircuit.Lupanov.any_comb}
\label{BoolCircuit.Lupanov.any_comb}
\leanok
\uses{BoolCircuit.Lupanov.assemble_heapPath, BoolCircuit.Lupanov.code, BoolCircuit.Lupanov.combRoots, BoolCircuit.Lupanov.gates, BoolCircuit.Lupanov.heapPath, BoolCircuit.Lupanov.heapPath_bounds, BoolCircuit.Lupanov.testBit_ofBits, BoolCircuit.Lupanov.value_comb, BoolCircuit.vertexValue}
\difficulty{5}
Let $n,k\in\mathbb{N}$ with $k\le n$, let $f:\{0,1\}^n\to\{0,1\}$, and let
$x\in\{0,1\}^n$. The Boolean disjunction of the values at all combination roots of the
circuit constructed from $f$ and $k$, when evaluated on input $x$, equals $f(x)$.
\end{theorem}

\begin{theorem}[Specification of the Lupanov gate construction]
\lean{BoolCircuit.Lupanov.lupGates_spec}
\label{BoolCircuit.Lupanov.lupGates_spec}
\leanok
\uses{BoolCircuit.Lupanov.combRoots, BoolCircuit.Lupanov.gates, BoolCircuit.Lupanov.gates_acyclic, BoolCircuit.Lupanov.length_gates, BoolCircuit.Lupanov.lupGates, BoolCircuit.Lupanov.offsets, BoolCircuit.TreeSpec, BoolCircuit.emitTree_spec}
\difficulty{4}
Let $n$ and $k$ be natural numbers with $k\le n$, and let $f$ be a Boolean function on
$n$-bit inputs. Start with the gate list prescribed by the Lupanov construction for $f$
and $k$, and take as roots the $2^{n-k}$ consecutive vertices beginning at the
construction's combining-gate block. The construction appends an acyclic list of OR
gates and returns an output vertex that lies within the resulting circuit; on every
input, its value is the disjunction of the values at those roots. The resulting list has
length at most the initial length plus the tree cost of the number of roots, and the
output depth is at most the ceiling of the base-$2$ logarithm of the number of roots,
plus one, plus the maximum depth of the roots in the initial circuit. Each appended gate
has at most two distinct inputs.
\end{theorem}

\begin{definition}[Lupanov circuit construction]
\lean{BoolCircuit.Lupanov.lupDAG}
\label{BoolCircuit.Lupanov.lupDAG}
\leanok
\uses{BoolCircuit.DAGCircuit, BoolCircuit.Lupanov.lupGates, BoolCircuit.Lupanov.lupGates_spec}
For each natural number $n$, Boolean function $f : \{0,1\}^n \to \{0,1\}$, and natural
number $k \le n$, the Lupanov construction assigns a directed acyclic circuit with $n$
input vertices. Its gates are ordered so that each gate reads only input vertices or
earlier gates, and its output is a vertex of the circuit.
\end{definition}

\begin{theorem}[Evaluation of the constructed circuit]
\lean{BoolCircuit.Lupanov.lupDAG_eval}
\label{BoolCircuit.Lupanov.lupDAG_eval}
\leanok
\uses{BoolCircuit.DAGCircuit.eval, BoolCircuit.Lupanov.any_comb, BoolCircuit.Lupanov.lupDAG, BoolCircuit.Lupanov.lupGates, BoolCircuit.Lupanov.lupGates_spec, BoolCircuit.vertexValue}
\difficulty{3}
Let $n$ be a natural number, let $f:\{0,1\}^n\to\{0,1\}$ be a Boolean function, and
let $k$ be a natural number with $k\le n$. The circuit constructed from $f$ and $k$
computes $f$: on every input $x\in\{0,1\}^n$, its output is $f(x)$.
\end{theorem}

\begin{theorem}[Fan-in bound for the constructed circuit]
\lean{BoolCircuit.Lupanov.lupDAG_isFaninTwo}
\label{BoolCircuit.Lupanov.lupDAG_isFaninTwo}
\leanok
\uses{BoolCircuit.DAGCircuit.IsFaninTwo, BoolCircuit.DAGGate.FaninTwo, BoolCircuit.Lupanov.gateAt_faninTwo, BoolCircuit.Lupanov.gates, BoolCircuit.Lupanov.lupDAG, BoolCircuit.Lupanov.lupGates, BoolCircuit.Lupanov.lupGates_spec}
\difficulty{5}
Let $n$ be a natural number, let $f$ be a Boolean function on $n$ bits, and let $k$ be a
natural number with $k\le n$. The circuit constructed from $f$ and $k$ is well formed,
and every gate reads at most two vertices.
\end{theorem}

\begin{theorem}[Size bound for the constructed circuit]
\lean{BoolCircuit.Lupanov.lupDAG_size_le}
\label{BoolCircuit.Lupanov.lupDAG_size_le}
\leanok
\uses{BoolCircuit.DAGCircuit.size, BoolCircuit.Lupanov.combRoots, BoolCircuit.Lupanov.length_gates, BoolCircuit.Lupanov.lupDAG, BoolCircuit.Lupanov.lupGates, BoolCircuit.Lupanov.lupGates_spec, BoolCircuit.Lupanov.offsets, BoolCircuit.treeCost}
\difficulty{4}
Let $n$ be a natural number, let $f$ be a Boolean function on $n$ bits, and let $k$ be a
natural number with $k\le n$. The size of the circuit constructed from $f$ and $k$ is
at most $2n+4\cdot 2^{n-k}+2\cdot 2^k+2^{2^k}$.
\end{theorem}

\begin{theorem}[Exponential bound for a square]
\lean{BoolCircuit.succ_sq_le_four_mul_two_pow}
\label{BoolCircuit.succ_sq_le_four_mul_two_pow}
\leanok
\difficulty{4}
For every natural number $n$, $(n+1)^2 \le 4\cdot 2^n$.
\end{theorem}

\begin{theorem}[Arithmetic bound for circuit size]
\lean{BoolCircuit.lupanov_arith}
\label{BoolCircuit.lupanov_arith}
\leanok
\uses{BoolCircuit.succ_sq_le_four_mul_two_pow}
\difficulty{5}
For every natural number $n$,
\[
(n+1)\bigl(2n+4\cdot 2^{\,n-(\lfloor\log_2 n\rfloor-1)}+2\cdot 2^{\,\lfloor\log_2
n\rfloor-1}+2^{\,2^{\lfloor\log_2 n\rfloor-1}}\bigr)\le 40\cdot 2^n,
\]
where the subtractions and logarithm are interpreted as natural-number operations, with
$\lfloor\log_2 n\rfloor$ denoting the base-two integer logarithm.
\end{theorem}

\begin{theorem}[Small fan-in-two circuit for every Boolean function]
\lean{BoolCircuit.exists_dagCircuit_faninTwo_mul_size_le}
\label{BoolCircuit.exists_dagCircuit_faninTwo_mul_size_le}
\leanok
\uses{BoolCircuit.DAGCircuit, BoolCircuit.DAGCircuit.IsFaninTwo, BoolCircuit.DAGCircuit.eval, BoolCircuit.DAGCircuit.size, BoolCircuit.Lupanov.lupDAG, BoolCircuit.Lupanov.lupDAG_eval, BoolCircuit.Lupanov.lupDAG_isFaninTwo, BoolCircuit.Lupanov.lupDAG_size_le, BoolCircuit.lupanov_arith}
\difficulty{3}
For every nonnegative integer $n$ and every Boolean function on assignments to $n$
variables, there is a directed acyclic circuit with those inputs and a valid output
vertex that computes the function on every assignment. Each gate has at most two inputs,
has no repeated input index, and every NOT gate has exactly one input; moreover, if the
circuit size counts the input vertices and gates, then $(n+1)$ times its size is at most
$40\cdot 2^n$.
\end{theorem}

\begin{theorem}[Small fan-in-two circuit representation]
\lean{BoolCircuit.exists_dagCircuit_faninTwo_size_le_div}
\label{BoolCircuit.exists_dagCircuit_faninTwo_size_le_div}
\leanok
\uses{BoolCircuit.DAGCircuit, BoolCircuit.DAGCircuit.IsFaninTwo, BoolCircuit.DAGCircuit.eval, BoolCircuit.DAGCircuit.size, BoolCircuit.exists_dagCircuit_faninTwo_mul_size_le}
\difficulty{3}
For every natural number $n$ and every Boolean function $f\colon\{0,1\}^n\to\{0,1\}$,
there is a well-formed acyclic Boolean circuit with at most two inputs to each gate that
computes $f$ and has size at most $40\cdot 2^n/(n+1)$.
\end{theorem}

\begin{theorem}[Universal circuit size bound]
\lean{Language.inSIZE_forty_two_pow_div}
\label{Language.inSIZE_forty_two_pow_div}
\leanok
\uses{BoolCircuit.DAGCircuit.IsFaninTwo, BoolCircuit.DAGCircuit.eval, BoolCircuit.DAGCircuit.size, BoolCircuit.DAGCircuitFamily, BoolCircuit.DAGCircuitFamily.mem_language_iff, BoolCircuit.exists_dagCircuit_faninTwo_size_le_div, Language.InSIZE}
\difficulty{3}
Every language $L$ over the Boolean alphabet is recognized by a family of circuits with
fan-in at most two whose circuit on inputs of length $n$ has size at most $40\cdot
2^n/(n+1)$ for every $n$.
\end{theorem}
